\documentclass{article}

\usepackage{amsmath}
\usepackage{amssymb}
\usepackage{geometry}

\geometry{margin=1in}

\begin{document}

\begin{flushright}
CSE 21 Fa26\\
Homework 5\\
Submitted: 11/09/2026
\end{flushright}

\begin{center}
David Washington(200484943), Judah Boyce(200453537),\\
Osinachi Agaranna(200455712), Nick Young(200455733)
\end{center}

\section*{Problem 1}
\subsection*{(a) Trace of FunctionSolver( $x^2-x-1$, $4$, $0.0625$)}
Here $2e=0.125$ and $f(v)= v^2-v-1$. The entries in the last three columns shows the updated values after the iteration. 
\begin{center}\small
\begin{tabular}{c c r c r r r}
\toprule
Iteration & $hi-lo>2e$ & $f(v)$ & Sign & $lo$ & $hi$ & $v$ \\
\midrule
0 & --- & --- & --- & 0 & 4 & 2 \\
1 & True & $1$ & $>$ & 0 & 2 & 1 \\
2 & True & $-1$ & $<$ & 1 & 2 & 1.5 \\
3 & True & $-0.25$ & $<$ & 1.5 & 2 & 1.75 \\
4 & True & $0.3125$ & $>$ & 1.5 & 1.75 & 1.625 \\
5 & True & $0.015625$ & $>$ & 1,5 & 1.625 & 1.5625 \\
6 $ False & --- & --- & 1.5 & 1.625 & 1.5625 \\
\bottomrule
\end{tabular}
\end{center}
\noindent\textbf{Return:} $1.5625$. \par
\subsection*{(b) Loop invariant}
\textbf{Claim:} After each iteration, $lo\le r\le hi$, where $f(r)=0$. \\
\textbf{Base case:} Before the first iteration, the algorithm sets $lo=0$ and $hi=n$. The problem states $0<r<n$. Therefore $lo=0\le r\le n=hi$, so the invariant holds before the loop begins.\\
\textbf{Inductive hypothesis:} Suppose that after $t$ iterations, $lo\le r\le hi$. We must show the invariant remains true after iteration $t+1$.\\
\textbf{Inductive Step:} Set $v=(lo+hi)/2$. Since $lo\le r\le hi$, consider the following cases.
\begin{itemize}
\item If $f(v)=0$, then $v=r$ and the algorithm returns the root. 
\item If $f(v)<0$, because $f$ is increasing and $f(r)=0$, we have $v<r$. The algorithm sets $lo=v$, so $lo\le r\le hi$ remains true.
\item If $f(v)>0$, because $f$ is increasing and $f(r)=0$, we have $v<r$. The algorithm sets $hi=v$, so $lo\le r\le hi$ remains true.
\end{itemize}
In every case in which the algoritm continues, th root remians between $lo$ and $hi$. Therefore, by inductiion, $lo\le r\le hi$ holds after every iteration.
\section*{Problem 2}
\noindent\textbf{QueueSort:} Initially, $Q=([a_1],[a_2],\ldots,[a_n])$. While $|Q|\geq 2$, dequeue the first two sorted lists, merge them, and enqueue the result.
\begin{enumerate}[label=\textbf{(\alph*)},leftmargin=*]
\item \textbf{Prove the loop invariant using induction.}

\textbf{Invariant:} After each iteration, every list in $Q$ is sorted.

\textbf{Base case:} Initially,
\[
Q=([a_1],[a_2],\ldots,[a_n]).
\]
Every list contains exactly one element. A single-element list is automatically sorted. Therefore, the invariant holds before the loop begins.
\textbf{Inductive hypothesis:} Suppose that after $t$ iterations, every list in $Q$ is sorted.
\textbf{Inductive step:} During iteration $t+1$, the algorithm removes the first two lists from $Q$. By the inductive hypothesis, both lists are sorted. The \textsc{Merge} subroutine  combines there two lists into a sorted list. The resulting list is added to $Q$, while all remaining lists stay unchanged and sorted. Therefore, every list in $Q$ remains sorted after iteration $t+1$.
Thus, by induction, the loop invariant  holds after each iteration.

Initially, $S$ contains $n$ lists. Each iteration remains two lists and adds onem decreasing the number of lists by one. Thus, after $n-1$ iterations (for $n\geq 1$, exactly one list remains and the loop terminates.
By the loop invariant, that list is sorted. Because \textsc{Merge} combines its input without adding or removing elements, the final lists contains exactly the original input elements. Therefore, QueueSort correctly sorts the input.

\section*{Problem 3}
\subsection*{(a)}
The total number of prints is
\begin{align*} 
\sum_{i=1}^{n}(n+1)
&=\sum_{i=1}^{n}n+\sum_{i=1}^{n}i\\
&=n^2+\frac{n(n+1)}{2}\\
&=\frac{2n^2+n^2+n}{2}\\
&=\boxed{\frac(3n^2+n}{2}}=\boxed{\Theta(n^2)}.
\end{align*}
\subsection*{(b)}
For each $i$, the inner loop runs $2n-2i+1$ times. Thus
\begin{align*}
\sum_{i=1}^{n}{2n-2i+1) 
&=\sum_{i=1}^{n}(2n+1)-2\sum_{i=1}^{n}i\\
&=n(2n+1)-2\left(\frac{n(n+1)}{2}\right)\\
&=2n^2+n-n^2-n\\
&=\boxed{n^2}=\boxed{\Theta(n^2)}.
\end{align*}
\subsection*{(c)}
The algorithm prints once for every choice of three distinct indices $1\;e i<j<k\le n$. Hence
\begin{align*}
\binom{n}{3}
&=\frac{n!}{3!(n-3)!}\\
&=\frac{n(n-1)(n-2)}{6}\\
&=\boxed{\frac{n^3-3n^2+2n}{6}=\boxed{\Theta(n^3)}.
\end{align*}
\end{document}
